\documentclass[11pt]{article}
\usepackage[letterpaper,margin=.68in,headheight=16pt,headsep=.18in,footskip=.28in]{geometry}
\usepackage{amsmath,array,enumitem,fancyhdr,hyperref}
\pagestyle{fancy}\fancyhf{}\renewcommand{\headrulewidth}{.3pt}
\fancyhead[L]{\small Bellingham Math Circle / Week 16 / Return visits}
\fancyhead[R]{\small Adult guide}
\fancyfoot[L]{\small RV-16-FAC-v1 / Draft and unpiloted}\fancyfoot[R]{\small\thepage}
\setlength{\emergencystretch}{2em}
\setlength{\parindent}{0pt}\setlength{\parskip}{7pt}
\setlist[itemize]{leftmargin=1.3em,itemsep=3pt,topsep=4pt}
\hypersetup{hidelinks,pdftitle={Week 16 return-visit facilitator guide}}
\begin{document}
{\Large\bfseries Mathematical overview}\par\medskip
\textbf{Facts and limits.} Count mesh cells: individual triangles with no mesh line inside them. Ordinary Sperner labeling uses only R/B/Y, fixes those three outer corners, and permits only the two endpoint labels along each side. Allowing G only at interior vertices defeats the guarantee of a specifically RBY cell. Nevertheless, if no RBY cell exists, at least three cells have three different labels: one each of RBG, RYG and BYG. Temporarily replace every G by R, B or Y in three separate copies of the unchanged original and apply ordinary Sperner; the forced rainbow in each copy identifies one of those three original types. On the three-step mesh the minimum three is attainable. All 256 legal four-label fillings include 27 with no RBY: 18 have three any-three cells and nine have five. These finite counts describe this mesh, not every triangulation.

For an edge-to-edge triangulated disk with arbitrary R/B/Y boundary labels, rainbow count has the parity of the number of boundary R-B edges. Interior doors count twice. A square can force even or odd counts; its boundary does not fix the count itself. Problems 3--4 use only R/B/Y, resetting Problem 1's four-label rule.

With outer corners R,B,Y counterclockwise and usual endpoint-label side restrictions, positive R-B-Y cells (up to cyclic reading) minus negative R-Y-B cells is exactly one. Oriented interior edges cancel, leaving the signed outer R-to-B changes. Reversing the outer corner orientation gives minus one. These statements require no holes or dangling subdivisions.

\textbf{Children's destination.} Change an interior-label assumption while preserving the boundary; distinguish RBY from any three different labels; investigate general boundary parity and opposite-sign pairs with one net positive.

\textbf{Entry map.} Problems 1--2 (pages 1--2), one investigation: oral K--1 star entry on page 1; grades 2--3 free mesh and separate counts on page 2; grades 4--5 minimum proof using restored copies. Problem 3 (page 3): grades 2--3 and up, boundary changes and even/odd pairing. Problem 4 (page 4): grades 4--5, a stable counterclockwise convention and two cyclic orders. Formal topology is unnecessary.

\textbf{New versus base.} Base K--1 Problem 4 already inserts centers and classifies local refinements; that is not a new investigation. These visits permit a fourth interior label, change the domain/boundary pattern, and count signs.
\newpage {\Large\bfseries Prepare a return visit}\par\medskip
\textbf{Status and use.} Three investigations extend one familiar theme. These companions are separate from the base packets and may support several visits. All pages are new and unpiloted. Printed labels describe prerequisites, not age restrictions. Offer one investigation’s pages at a time; completing every page is not the goal.

\textbf{Materials and preparation.} For each pair, four lettered/color pencils are the simplest preparation. For movable labels instead, supply sixteen each R/B/Y and two G; ruler, eraser and spare/tracing paper. Reserve six sets for eleven children. Counters should be at most 12 mm across with vertex spacing at least 18 mm. A three-step mesh has ten vertices, only one interior; all four square boards together have sixteen corners, which sets the movable-label maximum. Labels mark vertices, not faces. Save the original before recoloring and restore it before each replacement. G is allowed only at interior dots in Problems 1–2; remove it for Problems 3–4. Offer pages 1–2 together so page 1’s shared rules remain visible beside the free mesh. Printed dot rings are about 7.1 mm across; pencils or small labels can cover them, with counter centers aligned to the vertices. Material fit and orientation procedure are untested physically.

\textbf{Staffing.} Current context: eleven children, fixed KK11 / 3333 / 445 tables, one adult anchored to each table. Use pairs; the upper third child rotates as reader or checker. Routine adult reading, arrow drawing or recording can leave the mathematical decisions with children. If adults must choose every move, use the earlier entry rather than run the investigation for them.

\textbf{Short common launch.} Let children handle labels. Show how to label a dot and identify one cell’s three corner letters. Contrast a non-task RBG cell with an RBY cell: both have three different letters, only the second is specifically RBY. Use one count at a time. For signs trace counterclockwise around a separate R-Y-B triangle and show the intermediate read order and negative record. Turning that example preserves its sign. Do not reveal the minimum or signed theorem before exploration.

\textbf{Suggested first route.} Start with the small star under spoken rules, then offer the free mesh to children ready to choose boundary labels and preserve a filling. Seek zero RBY first; count any-three cells in a second pass. Save the recoloring proof for children who retain the unchanged reference. The square is a separate return direction. Introduce signs after rainbow recognition and a stable counterclockwise convention.

\textbf{Flexible timing.} Allow 3--5 minutes of material play and 2--4 minutes for the common demonstration, then 15--25 minutes for one investigation and 5 minutes to share. A longer second visit can spend 30--45 minutes on the same investigation, including unfinished examples or its explanation. These are planning estimates, not observed timings. Do not require all three within one hour.

\textbf{Questions and hints.} Start with ``What can you try?'' and ``What would count as success?'' Check the actual rule before giving a strategy. Optional questions and hints are keyed below; use them only after a real attempt. A counter argument or drawing may be a complete explanation.

\textbf{Source and pilot record.} Mathematical examples and arguments here are original local follow-ups to the current base theme. Consulted teaching source: \emph{Math Circle by the Bay}, Preface pp. vii--x (local PDF pp. 8--11), on interaction, manipulatives, hints, explanations and themes spanning multiple years. Our exact launch, entry route and timing are design inferences. Year-1 \emph{Handouts 6}, Problems 6.2--6.6, provided a local example of connecting representations. Record actual problem, starting route, examples tried, what needed adult rescue, and what children want to resume. Distinguish observations from hypotheses. Counter fit, materials stock, procedures and classroom response have not been physically tested.
\newpage
\textbf{Problems 1--2: a fourth interior label.} A star with R/B/Y outer corners and G at its center has zero RBY cells and three any-three cells, one each RBG, RYG and BYG. On the three-step mesh, read rows from the bottom R-to-B edge upwards: RRBB / RGB / YY / Y. Its nine cells comprise three monochrome, three using two labels, and one each RBG, RYG, BYG. Thus zero RBY and three any-three are attained on the larger mesh.

\textbf{Why three is the minimum.} Preserve the original no-RBY filling. In a copy replace every G by R. Ordinary Sperner forces an RBY cell. It was not originally RBY, so its original labels were G,B,Y. Restore the original and replace every G by B; this forces an original R,G,Y cell. Restore again and replace G by Y; this forces an original R,B,G cell. These three types are different, so at least three any-three cells exist. The general lower bound uses an argument, not numerical search.

\textbf{Conditions and limits.} The triangle is triangulated edge to edge, with usual R/B/Y corners and side restrictions and G only inside. A boundary G would invalidate the recolorings’ Sperner hypotheses. The star is an entry; the free mesh and minimum question supply the substantial investigation. Repeated star choices are not separate investigations.

\textbf{If needed.} Check one cell at a time and distinguish ``exactly RBY'' from ``any three different.'' Mark RBY first, then make a fresh pass for any-three. Adults may read and record while children choose. Three tracing copies or reversible cards support the proof without changing the original between tests.

\textbf{Deeper continuation.} Each of RBG, RYG and BYG occurs oddly if RBY is absent. For RB doors, RBG cells contribute one, RB-only cells two, and interior doors twice. The RB outer side has an odd number of changes. Repeat with RY and BY. The independent three-step enumeration checks 256 fillings: 27 have zero RBY, 18 of those have three any-three cells and nine have five.
\newpage
\textbf{Problem 3: a square boundary.} The student squares are unlabeled. For these adult examples number corners 1=lower left, 2=lower right, 3=upper right, 4=upper left. The left square has diagonal 1--3; the right has 2--4. For a square with diagonal 1--3, use cells 123 and 134; the other diagonal gives cells 124 and 234. With boundary RBRB there are four R-B doors and zero rainbows. With RBY Y there is one door and one rainbow under either diagonal. With RBRY there are two doors: diagonal 1--3 gives zero rainbows, while 2--4 gives two. Spaces in the boundary word only aid reading; the four corner order is fixed.

\textbf{General parity reason.} A rainbow cell has one R-B door. A cell using R and B only has two doors. Other cells have none. Count cell-door incidences: the rainbow count plus twice the two-door-cell count equals the boundary-door count plus twice the interior-door count. Therefore rainbow count and boundary-door count have the same odd/even parity. This works for any triangulated disk, even when the domain is a square or its boundary lacks the usual side restrictions.

A child can pair the two incidences of every inside door, leaving the outside doors unpaired. The remaining parity is easier to see physically than to state symbolically.

\textbf{If needed.} Trace only the boundary once and mark its R-B changes. Keep the chosen diagonal or refinement fixed while filling labels. Supply a pattern with two doors after children have discovered an odd case.

\textbf{Extension.} Find two fillings with the same boundary and different totals. RBRY on the two diagonal choices supplies zero and two; refining a two-label cell gives another construction. Same boundary fixes parity, not exact count.
\newpage
\textbf{Problem 4: positive and negative rainbows.} On the standard three-step triangle, all 192 legal fillings have signed counts \((1,0),(2,1)\) or \((3,2)\). The respective numbers of fillings are 108, 72 and 12. In each, positive minus negative is one. These finite counts are examples; the theorem applies to every edge-to-edge Sperner triangulation with the stated boundary orientation.

\textbf{Why the signed count is one.} Traverse each cell counterclockwise. Give an R-to-B edge contribution +1, a B-to-R edge -1, and every other edge zero. A positive rainbow contributes +1 and a negative one -1. A two-label R/B cell has one edge in each direction and contributes zero; every other nonrainbow also contributes zero.

When adding cell contributions, each interior edge occurs in both directions and cancels. Along the outer R-to-B side only R/B labels occur: +1 changes and -1 changes alternate, beginning with R and ending with B, so their net is +1. The other two sides have no R-B edge. Thus the whole boundary sum is +1, which equals positive count minus negative count.

\textbf{Conventions.} RBY, BYR and YRB are the same positive cyclic order; RYB, YBR and BRY are negative. Turning the page preserves counterclockwise orientation; reflecting it reverses signs. Fixed outer corners must be R,B,Y counterclockwise for the +1 statement. Reversing that order gives -1.

\textbf{If needed.} Put a small curved arrow inside each candidate rainbow and read around it aloud. Let the child classify one triangle at a time; adults may record the plus/minus marks. If the adult must operate every orientation, return to the four-label or square page.

\textbf{Extension.} Apply a missing-color refinement inside a two-label triangle. It creates one positive and one negative rainbow, directly illustrating preservation of the signed difference.
\end{document}
